% ats-single-column — a plain, single-column, ATS-safe CV.
%
% Compiles with Tectonic using bundled resources only: the font is TeX Gyre
% Heros, loaded by file name so no system font install is required.
%
% The renderer substitutes %%PDFMETA%% and %%BODY%%. Everything else is design.
% The macro contract this template must provide is in PLACEHOLDERS.md.

\documentclass[11pt,a4paper]{article}

\usepackage[a4paper,top=1.1cm,bottom=1.1cm,left=1.4cm,right=1.4cm]{geometry}
\usepackage{fontspec}
\setmainfont{texgyreheros}[
  Extension      = .otf,
  UprightFont    = *-regular,
  BoldFont       = *-bold,
  ItalicFont     = *-italic,
  BoldItalicFont = *-bolditalic,
]
\usepackage{xcolor}
\usepackage{enumitem}
\usepackage{titlesec}
\usepackage{microtype}
\usepackage[normalem]{ulem}
\usepackage{hyperref}

\definecolor{rulegray}{gray}{0.60}
\definecolor{muted}{gray}{0.30}
\definecolor{linkblue}{HTML}{0B57D0}

% Links are blue and underlined: on screen that is the only signal that a line
% of text is clickable, and a printed copy loses the colour but keeps the rule.
\hypersetup{
  colorlinks = true,
  urlcolor   = linkblue,
  linkcolor  = linkblue,
  breaklinks = true,
}

% The renderer emits plain \href, so the underline is applied here, where the
% design lives. \uline rather than \underline because it lets a long URL break
% across a line instead of running into the margin.
\let\cvhref\href
\renewcommand{\href}[2]{\cvhref{#1}{\uline{#2}}}

\pagestyle{empty}
\setlength{\parindent}{0pt}
\setlength{\parskip}{0pt}
\linespread{1.0}
\raggedbottom

% Ragged right, not justified: a CV column is narrow enough that justification
% forces mid-word hyphenation ("main-tainability"), which reads badly and can
% split a keyword an ATS is looking for. \sloppy keeps lines from overflowing
% the margin now that they cannot be stretched.
\raggedright
\frenchspacing
\hyphenpenalty=10000
\exhyphenpenalty=10000
\sloppy

% --- section heading: bold, sentence case, thin rule underneath -------------
\titleformat{\section}
  {\normalfont\large\bfseries}{}{0em}{}
  [\vspace{-0.55em}{\color{rulegray}\rule{\linewidth}{0.5pt}}]
\titlespacing*{\section}{0pt}{0.55em}{0.25em}

% ===========================================================================
% Macro contract consumed by render_cv.py — see PLACEHOLDERS.md
% ===========================================================================

% Header. \centering inside a group rather than the center environment, which
% would add \topsep above and below and loosen the block the name sits in.
\newcommand{\cvname}[1]{{\centering\fontsize{18}{20}\selectfont\bfseries #1\par}\vspace{0.15em}}
\newcommand{\cvheadline}[1]{{\centering\large #1\par}\vspace{0.2em}}
\newcommand{\cvcontact}[1]{{\centering\small\color{muted}#1\par}\vspace{0.15em}}

% Sections
\newcommand{\cvsection}[1]{\section*{#1}}
\newcommand{\cvsummary}[1]{#1\par\vspace{0.1em}}

% \raggedright leaves \rightskip at 0pt plus 1fil, so a plain \hfill would
% share the slack with it and stop halfway across the line. A second-order
% fill outranks first-order fil and takes the whole gap.
\newcommand{\cvfill}{\hskip 0pt plus 1fill\relax}

% \cvitem{headline}{metadata}
%
% Dates and location sit at the right margin on the SAME line as the headline.
% That is one paragraph, not a second column, and the distinction is the whole
% point: a real column (minipage, tabular, paracol) is emitted as its own text
% block, and PDF extraction then reattaches a date to the wrong entry. Inside
% one line box the extractor reads "headline ... metadata" in order.
% build_and_validate.py re-checks that on every build.
%
% The metadata is boxed, so it is this or nothing: where the pair does not fit
% on one line TeX breaks at \cvfill and the whole of it drops to the next line
% flush left — the old layout, still correct. Left unboxed it would instead
% break wherever it ran out of room, stranding half a date range on a line of
% its own. \footnotesize buys the width that keeps the longest rows out of
% that fallback in the first place.
\newcommand{\cvitem}[2]{%
  \vspace{0.3em}%
  #1\cvfill\mbox{\footnotesize\color{muted}#2}\par
}

% \cvitemplain{headline} — same, for entries with no metadata line
\newcommand{\cvitemplain}[1]{\vspace{0.3em}#1\par}

% \cvsubline{text} — ONE continuation line under an entry headline, for a
% fact that belongs to the entry but not on its headline (an institution and
% its grade).
%
% One is the limit, and it is a measured one. With \cvitem's metadata at the
% right margin, two or more flush-left lines beneath it turn the entry into a
% block poppler reads as two columns: it emits the whole left column first,
% and the metadata surfaces later — under the next section heading, or worse,
% on the next entry, where nothing marks it as the wrong date. Anything extra
% belongs on this line, not on another one.
\newcommand{\cvsubline}[1]{#1\par}

% Italic metadata line (tech stack, etc.)
\newcommand{\cvmeta}[1]{{\small\itshape #1}\par}

% One project URL, already labelled by the renderer ("Demo: ...", "Git: ...").
% Emitted once per link, so each gets its own line.
\newcommand{\cvlinks}[1]{{\small\color{muted}#1}\par}

% Achievement bullets
\newlist{cvbullets}{itemize}{1}
\setlist[cvbullets]{
  leftmargin = 1.15em,
  labelsep   = 0.45em,
  itemsep    = 0.05em,
  parsep     = 0pt,
  topsep     = 0.12em,
  partopsep  = 0pt,
  label      = \textbullet,
}

% Skills line: {group label}{comma-separated items}. Also carries a language
% row, which has the same shape: a bold label and one value.
\newcommand{\cvskill}[2]{\vspace{0.08em}{\bfseries #1:} #2\par}

% Narrow-column forms, part of the macro contract. This template has no
% sidebar, so the renderer never emits them here.
\newcommand{\cvskilllabel}[1]{\vspace{0.2em}{\bfseries #1}\par}
\newcommand{\cvskillitem}[1]{{\small #1}\par}

% Single-line row, flush left (certifications)
\newcommand{\cvplain}[1]{\vspace{0.1em}#1\par}

\hypersetup{
  pdftitle={Nguyen Thi Dieu Hang — CV — Junior Frontend Engineer},
  pdfauthor={Nguyen Thi Dieu Hang},
  pdfcreator={},
}

\begin{document}
\cvname{Nguyen Thi Dieu Hang}
\cvheadline{Frontend Developer}
\cvcontact{\href{mailto:ntdieuhang192@gmail.com}{ntdieuhang192@gmail.com} $\cdot$ +84 327 840 518 $\cdot$ \href{https://www.linkedin.com/in/hangnguyen118}{linkedin.com/in/hangnguyen118} $\cdot$ \href{https://github.com/hangnguyen118}{github.com/hangnguyen118} $\cdot$ \href{https://browsermind-profile.vercel.app}{browsermind-profile.vercel.app}}
\cvsection{Summary}
\cvsummary{Frontend Developer and Information Technology graduate from Ho Chi Minh City University of Technology (HUTECH) with hands-on experience building responsive web and mobile applications using Next.js, React.js, and React Native.}
\cvsection{Education}
\cvitem{\textbf{Engineer's Degree, Information Technology}}{Jan 2025 $\cdot$ Ho Chi Minh City, Vietnam}
\cvsubline{Ho Chi Minh City University of Technology (HUTECH) $\cdot$ GPA 3.35 / 4.0 $\cdot$ Very Good}
\cvsection{Technical Skills}
\cvskill{Frontend}{React.js, TypeScript, JavaScript, HTML5, CSS3, Responsive design}
\cvskill{AI workflows}{Claude / Claude Code, Spec-Driven Development, Playwright MCP}
\cvskill{Data \& delivery}{TanStack Query, REST API integration, Git, Vercel}
\cvsection{Certifications}
\cvplain{Microsoft Certified: Azure Fundamentals (AZ-900) $\cdot$ \href{https://learn.microsoft.com/api/credentials/share/vi-vn/hangnguyen1108/306DA7D144E07FF7?sharingId=C357452C5EB8224A}{Verify Credential}}
\cvsection{Languages}
\cvskill{English}{Intermediate (Effective professional communication)}
\cvskill{Vietnamese}{Native}
\cvsection{Experience}
\cvitem{\textbf{Frontend Developer}, VILIHA PTE. LTD.}{Ho Chi Minh City, Vietnam $\cdot$ Jul 2026 – Present $\cdot$ Full-time}
\begin{cvbullets}
  \item Translate detailed Figma designs into responsive, semantic frontend code and reusable components.
  \item Develop with Next.js App Router, Server Components, TanStack Query, Tailwind CSS, and shadcn/ui.
  \item Collaborate in Agile/Scrum and communicate in English with project stakeholders.
\end{cvbullets}
\cvitem{\textbf{Frontend Developer \& Testing}, Chip Nova}{Ho Chi Minh City, Vietnam $\cdot$ Aug 2025 – Mar 2026}
\begin{cvbullets}
  \item Clarified requirements, designed user experiences, and shipped frontend features across web and mobile.
  \item Wrote Playwright end-to-end tests to verify product quality.
\end{cvbullets}
\cvitem{\textbf{Web Developer Intern}, Viet Japan Partner Company}{Ho Chi Minh City, Vietnam $\cdot$ May 2024 – Jun 2024 $\cdot$ Internship}
\begin{cvbullets}
  \item Developed web interfaces, integrated APIs, and collaborated on feature development and testing.
\end{cvbullets}
\cvsection{Projects}
\cvitem{\textbf{Portfolio Website (BrowserMind)}}{Frontend Developer \& Testing $\cdot$ Jun 2026}
\cvlinks{Demo: \href{https://browsermind-profile.vercel.app}{https://browsermind-profile.vercel.app}}
\cvlinks{Git: \href{https://github.com/hangnguyen118/browsermind-profile}{https://github.com/hangnguyen118/browsermind-profile}}
\begin{cvbullets}
  \item Built a React/TypeScript portfolio with Spec-Driven Development and Claude Code; captured requirements in SPEC.md and developed 23 React components.
  \item Used AI-powered development and testing workflows with Playwright MCP for automated validation.
\end{cvbullets}
\cvitem{\textbf{Shopee UI Test Automation}}{Test Automation Engineer $\cdot$ Jul 2026}
\cvlinks{Git: \href{https://github.com/hangnguyen118/agentic-test-framework-playwright-ai}{https://github.com/hangnguyen118/agentic-test-framework-playwright-ai}}
\begin{cvbullets}
  \item Built a Playwright/TypeScript framework with Page Object Model: 21 test cases across 4 browser and device targets, with traces and screenshots for debugging.
\end{cvbullets}

\end{document}
